\chapter{Boolean Analysis — Restriction Fourier}
%%% generated by the blueprint pipeline from TCSlib.BooleanAnalysis.LMN.RestrictionFourier
\section{Overview}

This module analyses how a random restriction acts on the Fourier spectrum of a
real-valued Boolean function $f : \{0,1\}^n \to \bbr$. It gives a closed form for the
Fourier coefficients of the restricted function $f_\rho$, factors the Bernoulli
restriction measure coordinate by coordinate, and derives O'Donnell's Proposition 4.17:
$\E_\rho[\widehat{f_\rho}(S)] = p^{\abs{S}}\widehat{f}(S)$ and
$\E_\rho[\widehat{f_\rho}(S)^2] = \sum_{U \supseteq S} p^{\abs{S}}(1-p)^{\abs{U \setminus S}}\widehat{f}(U)^2$,
together with the probability form $\Pr_\rho[U \cap J = S] = p^{\abs{S}}(1-p)^{\abs{U\setminus S}}$
for the set $J$ of free coordinates.

%%%%%%%%%%%%%%%%%%%%%%%%%%%%%%%%%%%%%%%%%%%%%%%%%%%%%%%%%%%%%%%%%%%%%%%%%%%%%%%
\section{Declarations}

\begin{definition}[Restriction commutes with the sign encoding]
\lean{RestrictionFourier.restrictBF_boolToSign}
\label{RestrictionFourier.restrictBF_boolToSign}
\leanok
\uses{BooleanAnalysis.boolToSign, SwitchingLemma2.Restriction, SwitchingLemma2.restrictFn}
\difficulty{1}
Write $\sigma\colon\{0,1\}\to\{-1,1\}$ for the sign encoding, $\sigma(0)=1$ and
$\sigma(1)=-1$, and for a function $g$ on the cube (Boolean- or real-valued) let
$g_\rho$ denote its restriction along $\rho$, obtained by reading each free coordinate
from the input and each fixed coordinate from $\rho$. Let $f\colon\{0,1\}^n\to\{0,1\}$
be a Boolean-valued function and let $\rho$ be a restriction on $n$ variables. Then
restricting the real-valued function $x\mapsto\sigma(f(x))$ along $\rho$ agrees with
first restricting $f$ and then applying the sign encoding; that is, $(\sigma\circ
f)_\rho=\sigma\circ f_\rho$, so that for every point $x$,
\[
  (\sigma\circ f)_\rho(x)\;=\;\sigma\big(f_\rho(x)\big).
\]
\end{definition}

\begin{definition}[Characterizing the free coordinates of a restriction]
\lean{RestrictionFourier.mem_freeVars}
\label{RestrictionFourier.mem_freeVars}
\leanok
\uses{SwitchingLemma2.Restriction, SwitchingLemma2.Restriction.freeVars}
\difficulty{1}
Let $\rho : \mathrm{Fin}\,n \to \mathrm{Option}\,\mathrm{Bool}$ be a restriction on $n$
variables, and let $i \in \mathrm{Fin}\,n$ be a coordinate. Then $i$ belongs to the set
of free variables of $\rho$ if and only if $\rho(i)$ is unset; that is, exactly when
$\rho$ leaves the coordinate $i$ free rather than fixing it to a Boolean value.
\end{definition}

\begin{lemma}[Splitting a Walsh character along a restriction]
\lean{RestrictionFourier.chiS_extend}
\label{RestrictionFourier.chiS_extend}
\leanok
\uses{BooleanAnalysis.BoolCube, BooleanAnalysis.chiS, RestrictionFourier.mem_freeVars, SwitchingLemma2.Restriction, SwitchingLemma2.Restriction.extend, SwitchingLemma2.Restriction.freeVars}
\difficulty{3}
Let $\rho$ be a restriction on $n$ variables, let $J\subseteq[n]$ be its set of free
coordinates, and let $U\subseteq[n]$ and $x\in\{0,1\}^n$ be arbitrary. Write $\hat x$
for the extension of $x$ by $\rho$, namely the point of $\{0,1\}^n$ that agrees with $x$
on the free coordinates and takes the fixed value of $\rho$ on the remaining ones. Then
the Walsh character $\chi_U$ evaluated at $\hat x$ factors as
\[
  \chi_U(\hat x) \;=\; \chi_{U\cap J}(x)\cdot\prod_{i\in U\setminus J}\sigma(b_i),
\]
where $b_i$ denotes the Boolean value that $\rho$ fixes at coordinate $i$ and $\sigma$
is the sign encoding sending $\mathtt{false}\mapsto 1$ and $\mathtt{true}\mapsto -1$.
Thus the character splits into a Walsh character in the free coordinates lying in $U$
times a constant $\pm1$ sign contributed by the fixed coordinates of $U$.
\end{lemma}

\begin{theorem}[Fourier coefficients of a restricted Boolean function]
\lean{RestrictionFourier.fourierCoeff_restrictBF}
\proofsource{boolean-ch03-spectral-learning-restrictions}{
  pdf-20e15a3d3d5e-p008-b002,
  pdf-20e15a3d3d5e-p008-b003
}
\label{RestrictionFourier.fourierCoeff_restrictBF}
\leanok
\statementsource{boolean-ch03-spectral-learning-restrictions}{
  pdf-20e15a3d3d5e-p008-b002,
  pdf-20e15a3d3d5e-p008-b003
}
\proofstep
  {BoolCircuit.Lit}
  {context}
  {boolean-ch03-spectral-learning-restrictions}
  {pdf-20e15a3d3d5e-p005-b002}
\proofstep
  {BoolCircuit.Lit.eval}
  {context}
  {boolean-ch03-spectral-learning-restrictions}
  {pdf-20e15a3d3d5e-p005-b002}
\proofstep
  {Literal}
  {context}
  {boolean-ch03-spectral-learning-restrictions}
  {pdf-20e15a3d3d5e-p005-b002}
\proofstep
  {DecisionTree.eval}
  {context}
  {boolean-ch03-spectral-learning-restrictions}
  {pdf-20e15a3d3d5e-p005-b002}
\proofstep
  {DecisionTree}
  {direct}
  {boolean-ch03-spectral-learning-restrictions}
  {pdf-20e15a3d3d5e-p005-b002}
\proofstep
  {DecisionTree.depth}
  {direct}
  {boolean-ch03-spectral-learning-restrictions}
  {pdf-20e15a3d3d5e-p006-b001}
\proofstep
  {SwitchingLemma2.Restriction}
  {context}
  {boolean-ch03-spectral-learning-restrictions}
  {pdf-20e15a3d3d5e-p007-b001,
    pdf-20e15a3d3d5e-p006-b005}
\proofstep
  {SwitchingLemma2.Restriction.freeVars}
  {context}
  {boolean-ch03-spectral-learning-restrictions}
  {pdf-20e15a3d3d5e-p007-b001}
\proofstep
  {SwitchingLemma2.Restriction.extend}
  {context}
  {boolean-ch03-spectral-learning-restrictions}
  {pdf-20e15a3d3d5e-p007-b001}
\proofstep
  {SwitchingLemma2.Literal.fixedBy}
  {context}
  {boolean-ch03-spectral-learning-restrictions}
  {pdf-20e15a3d3d5e-p007-b001}
\proofstep
  {Literal.flipNeg}
  {context}
  {boolean-ch03-spectral-learning-restrictions}
  {pdf-20e15a3d3d5e-p005-b002}
\proofstep
  {Literal.flipNeg_eval}
  {context}
  {boolean-ch03-spectral-learning-restrictions}
  {pdf-20e15a3d3d5e-p005-b002}
\proofstep
  {Literal.flipNeg_var}
  {context}
  {boolean-ch03-spectral-learning-restrictions}
  {pdf-20e15a3d3d5e-p005-b002}
\proofstep
  {DecisionTree.negateLeaves}
  {context}
  {boolean-ch03-spectral-learning-restrictions}
  {pdf-20e15a3d3d5e-p005-b002}
\proofstep
  {DecisionTree.negateLeaves_eval}
  {context}
  {boolean-ch03-spectral-learning-restrictions}
  {pdf-20e15a3d3d5e-p005-b002}
\proofstep
  {DecisionTree.negateLeaves_depth}
  {context}
  {boolean-ch03-spectral-learning-restrictions}
  {pdf-20e15a3d3d5e-p006-b001}
\proofstep
  {LMN.clauseIsTaut}
  {context}
  {boolean-ch03-spectral-learning-restrictions}
  {pdf-20e15a3d3d5e-p005-b002}
\proofstep
  {LMN.instDecidableClauseIsTaut}
  {context}
  {boolean-ch03-spectral-learning-restrictions}
  {pdf-20e15a3d3d5e-p005-b002}
\proofstep
  {LMN.mergeGates}
  {context}
  {boolean-ch03-spectral-learning-restrictions}
  {pdf-20e15a3d3d5e-p005-b002}
\proofstep
  {LMN.mergeGates_castAdd}
  {context}
  {boolean-ch03-spectral-learning-restrictions}
  {pdf-20e15a3d3d5e-p005-b002}
\proofstep
  {LMN.mergeGates_natAdd}
  {context}
  {boolean-ch03-spectral-learning-restrictions}
  {pdf-20e15a3d3d5e-p005-b002}
\proofstep
  {BooleanAnalysis.BoolCube}
  {context}
  {boolean-ch03-spectral-learning-restrictions}
  {pdf-20e15a3d3d5e-p001-b003}
\proofstep
  {BooleanAnalysis.BooleanFunc}
  {context}
  {boolean-ch03-spectral-learning-restrictions}
  {pdf-20e15a3d3d5e-p001-b003}
\proofstep
  {BooleanAnalysis.uniformWeight}
  {context}
  {boolean-ch03-spectral-learning-restrictions}
  {pdf-20e15a3d3d5e-p012-b002}
\proofstep
  {BooleanAnalysis.expect}
  {context}
  {boolean-ch03-spectral-learning-restrictions}
  {pdf-20e15a3d3d5e-p012-b002}
\proofstep
  {BooleanAnalysis.innerProduct}
  {context}
  {boolean-ch03-spectral-learning-restrictions}
  {pdf-20e15a3d3d5e-p012-b002}
\proofstep
  {BooleanAnalysis.boolToSign}
  {context}
  {boolean-ch03-spectral-learning-restrictions}
  {pdf-20e15a3d3d5e-p003-b005}
\proofstep
  {BooleanAnalysis.boolToSign_false}
  {context}
  {boolean-ch03-spectral-learning-restrictions}
  {pdf-20e15a3d3d5e-p003-b005}
\proofstep
  {BooleanAnalysis.boolToSign_true}
  {context}
  {boolean-ch03-spectral-learning-restrictions}
  {pdf-20e15a3d3d5e-p003-b005}
\proofstep
  {BooleanAnalysis.boolToSign_sq}
  {context}
  {boolean-ch03-spectral-learning-restrictions}
  {pdf-20e15a3d3d5e-p003-b005}
\proofstep
  {BooleanAnalysis.boolToSign_mul_self}
  {context}
  {boolean-ch03-spectral-learning-restrictions}
  {pdf-20e15a3d3d5e-p003-b005}
\proofstep
  {BooleanAnalysis.chiS}
  {direct}
  {boolean-ch03-spectral-learning-restrictions}
  {pdf-20e15a3d3d5e-p003-b005}
\proofstep
  {BooleanAnalysis.chiS_empty}
  {context}
  {boolean-ch03-spectral-learning-restrictions}
  {pdf-20e15a3d3d5e-p003-b005}
\proofstep
  {BooleanAnalysis.chiS_singleton}
  {context}
  {boolean-ch03-spectral-learning-restrictions}
  {pdf-20e15a3d3d5e-p003-b005}
\proofstep
  {BooleanAnalysis.chiS_mul_chiS}
  {direct}
  {boolean-ch03-spectral-learning-restrictions}
  {pdf-20e15a3d3d5e-p003-b005}
\proofstep
  {BooleanAnalysis.fourierCoeff}
  {direct}
  {boolean-ch03-spectral-learning-restrictions}
  {pdf-20e15a3d3d5e-p012-b002}
\proofstep
  {BooleanAnalysis.sum_prod_subset_eq_prod_one_add}
  {context}
  {boolean-ch03-spectral-learning-restrictions}
  {pdf-20e15a3d3d5e-p004-b005}
\proofstep
  {BooleanAnalysis.sum_chiS_mul_eq}
  {context}
  {boolean-ch03-spectral-learning-restrictions}
  {pdf-20e15a3d3d5e-p003-b005}
\proofstep
  {BooleanAnalysis.walsh_expansion}
  {direct}
  {boolean-ch03-spectral-learning-restrictions}
  {pdf-20e15a3d3d5e-p003-b005}
\proofstep
  {BooleanAnalysis.sum_boolToSign}
  {context}
  {boolean-ch03-spectral-learning-restrictions}
  {pdf-20e15a3d3d5e-p003-b005}
\proofstep
  {BooleanAnalysis.sum_chiS}
  {context}
  {boolean-ch03-spectral-learning-restrictions}
  {pdf-20e15a3d3d5e-p004-b001}
\proofstep
  {BooleanAnalysis.fourier_coeff_chi}
  {context}
  {boolean-ch03-spectral-learning-restrictions}
  {pdf-20e15a3d3d5e-p004-b001}
\proofstep
  {BooleanAnalysis.innerProduct_chi_self}
  {context}
  {boolean-ch03-spectral-learning-restrictions}
  {pdf-20e15a3d3d5e-p004-b001}
\proofstep
  {BooleanAnalysis.flipBit}
  {context}
  {boolean-ch03-spectral-learning-restrictions}
  {pdf-20e15a3d3d5e-p003-b003}
\proofstep
  {BooleanAnalysis.flipBit_flipBit}
  {context}
  {boolean-ch03-spectral-learning-restrictions}
  {pdf-20e15a3d3d5e-p003-b003}
\proofstep
  {BooleanAnalysis.instAddCommGroupBooleanFunc}
  {context}
  {boolean-ch03-spectral-learning-restrictions}
  {pdf-20e15a3d3d5e-p004-b001}
\proofstep
  {BooleanAnalysis.instModuleRealBooleanFunc}
  {context}
  {boolean-ch03-spectral-learning-restrictions}
  {pdf-20e15a3d3d5e-p004-b001}
\proofstep
  {BooleanAnalysis.boolToSign_not}
  {context}
  {boolean-ch03-spectral-learning-restrictions}
  {pdf-20e15a3d3d5e-p003-b005}
\proofstep
  {DecisionTree.fourierCoeff_sum_chiS}
  {context}
  {boolean-ch03-spectral-learning-restrictions}
  {pdf-20e15a3d3d5e-p003-b005}
\proofstep
  {RestrictionFourier.restrictBF}
  {context}
  {boolean-ch03-spectral-learning-restrictions}
  {pdf-20e15a3d3d5e-p007-b001}
\proofstep
  {RestrictionFourier.signProd}
  {context}
  {boolean-ch03-spectral-learning-restrictions}
  {pdf-20e15a3d3d5e-p008-b002}
\proofstep
  {RestrictionFourier.mem_freeVars}
  {context}
  {boolean-ch03-spectral-learning-restrictions}
  {pdf-20e15a3d3d5e-p007-b001}
\proofstep
  {RestrictionFourier.chiS_extend}
  {context}
  {boolean-ch03-spectral-learning-restrictions}
  {pdf-20e15a3d3d5e-p008-b003}
\uses{BooleanAnalysis.BooleanFunc, BooleanAnalysis.chiS, BooleanAnalysis.fourierCoeff, BooleanAnalysis.walsh_expansion, DecisionTree.fourierCoeff_sum_chiS, RestrictionFourier.chiS_extend, SwitchingLemma2.Restriction, SwitchingLemma2.Restriction.extend, SwitchingLemma2.Restriction.freeVars}
\difficulty{5}
Let $f : \{0,1\}^n \to \bbr$, and let $\rho$ be a restriction on $n$ variables; write $J
\subseteq [n]$ for its set of free coordinates, namely the coordinates left unset by
$\rho$. Let $f_\rho : \{0,1\}^n \to \bbr$ be the restricted function whose value at a
point $x$ is obtained by evaluating $f$ at the assignment that agrees with $x$ on the
free coordinates in $J$ and takes the bit fixed by $\rho$ on every other coordinate.
Then for every $S \subseteq [n]$,
\[
\widehat{f_\rho}(S) \;=\; \sum_{\substack{U \subseteq [n] \\ U \cap J = S}}
\widehat{f}(U)\;\prod_{i \in U \setminus J} \sigma(\rho_i),
\]
where $\sigma$ denotes the sign encoding and, for each $i \in U \setminus J$ (which lies
outside $J$ and is therefore a coordinate fixed by $\rho$), $\rho_i \in
\{\mathtt{false},\mathtt{true}\}$ is the bit that $\rho$ assigns there. This is
O'Donnell's Corollary 3.22.
\end{theorem}

\begin{lemma}[Sums of coordinatewise products over restrictions factor]
\lean{RestrictionFourier.sum_restriction_prod}
\label{RestrictionFourier.sum_restriction_prod}
\leanok
\uses{SwitchingLemma2.Restriction}
\difficulty{2}
Fix $n \in \bbn$, and for each coordinate $i \in \{1,\dots,n\}$ let $h_i :
\{\mathtt{true},\mathtt{false},\ast\} \to \bbr$ be a real-valued weight assigned to the
three possible states of that coordinate (fixed to true, fixed to false, or left free).
Then summing the product of these weights over all $3^n$ restrictions $\rho$ on $n$
variables equals the product over coordinates of the coordinatewise total weights:
\[
\sum_{\rho} \prod_{i=1}^{n} h_i\bigl(\rho(i)\bigr) \;=\; \prod_{i=1}^{n} \sum_{v}
h_i(v),
\]
where $\rho$ ranges over all restrictions on $n$ variables and $v$ ranges over the three
per-coordinate states.
\end{lemma}

\begin{lemma}[Bernoulli restriction weights factor over coordinates]
\lean{RestrictionFourier.sum_bernoulli_prod}
\label{RestrictionFourier.sum_bernoulli_prod}
\leanok
\uses{LMN.bernoulliRestrWeight_eq_prod, LMN.varWeight, RestrictionFourier.sum_restriction_prod, SwitchingLemma2.Restriction, SwitchingLemma2.bernoulliRestrWeight}
\difficulty{3}
Fix $n \in \bbn$ and a real parameter $p$, and recall that a restriction on $n$
variables is a map $\rho : \mathrm{Fin}\,n \to \mathrm{Option}\,\mathrm{Bool}$ assigning
each variable either a fixed Boolean value or the free marker $\mathrm{none}$. Write
$w_p(\rho) = p^{\,k}\bigl(\tfrac{1-p}{2}\bigr)^{\,n-k}$ for its Bernoulli weight, where
$k$ is the number of free variables of $\rho$, and let $\omega_p$ be the per-variable
weight given by $\omega_p(\mathrm{none}) = p$ and $\omega_p(v) = \tfrac{1-p}{2}$ for a
fixed value $v$. Then for every family $h_i : \mathrm{Option}\,\mathrm{Bool} \to \bbr$
indexed by $i \in \mathrm{Fin}\,n$,
\[
  \sum_{\rho}\, w_p(\rho) \prod_{i} h_i(\rho(i))
  \;=\; \prod_{i} \sum_{v}\, \omega_p(v)\, h_i(v),
\]
where the outer sum ranges over all restrictions $\rho$ on $n$ variables and each inner
sum ranges over $v \in \mathrm{Option}\,\mathrm{Bool}$.
\end{lemma}

\begin{definition}[Per-coordinate local factor]
\lean{RestrictionFourier.localFactor}
\label{RestrictionFourier.localFactor}
\leanok
\uses{BooleanAnalysis.boolToSign}
Given sets $U, S \subseteq [n]$, a coordinate $i$ and a per-coordinate value $v$, the
local factor is $1$ if $i \in S$ and $v = \mathtt{none}$ (the coordinate must be free),
$0$ if $i \in S$ and $v$ is fixed; for $i \in U \setminus S$ it is $0$ when $v =
\mathtt{none}$ and $\mathrm{boolToSign}(b)$ when $v = \mathtt{some}\ b$; and $1$ for all
remaining coordinates. It encodes the event $U \cap \mathrm{freeVars} = S$ together with
the sign contributed on $U \setminus S$.
\end{definition}

\begin{lemma}[Indicator times sign product factors over coordinates]
\lean{RestrictionFourier.indicator_signProd_eq_prod}
\label{RestrictionFourier.indicator_signProd_eq_prod}
\leanok
\uses{BooleanAnalysis.boolToSign, RestrictionFourier.localFactor, RestrictionFourier.mem_freeVars, SwitchingLemma2.Restriction, SwitchingLemma2.Restriction.freeVars}
\difficulty{5}
Let $U, S \subseteq [n]$ with $S \subseteq U$, and let $\rho$ be a restriction on $n$
variables whose free-variable set is $J$. Write $\sigma(b) \in \{\pm1\}$ for the sign
encoding of a bit $b$, and let $\ell(U,S,i,\rho(i))$ denote the per-coordinate local
factor at coordinate $i$. Then
\[
\mathbf{1}[\,U \cap J = S\,]\cdot \prod_{i \in U \setminus J} \sigma\bigl(\rho(i)\bigr)
  \;=\; \prod_{i=1}^{n} \ell\bigl(U, S, i, \rho(i)\bigr),
\]
where the left-hand side is read as $0$ whenever $U \cap J \neq S$. Equivalently, the
product on the right vanishes unless the free coordinates of $\rho$ lying in $U$ are
exactly $S$, in which case it equals the product of the $\pm1$-encodings of the bits
that $\rho$ fixes on $U \setminus J$.
\end{lemma}

\begin{lemma}[Averaging a single local factor against the Bernoulli weight]
\lean{RestrictionFourier.sum_varWeight_localFactor}
\label{RestrictionFourier.sum_varWeight_localFactor}
\leanok
\uses{BooleanAnalysis.boolToSign, LMN.varWeight, RestrictionFourier.localFactor}
\difficulty{2}
Fix a real parameter $p$, two coordinate sets $U, S \subseteq [n]$, and a coordinate $i
\in [n]$. Writing $w_p$ for the per-variable Bernoulli weight — which assigns $p$ to the
free value and $(1-p)/2$ to each of the two fixed values — and $\ell_{U,S,i}$ for the
per-coordinate local factor at $i$, the weighted sum over the three possible
per-coordinate values $v$ (free, or fixed to either Boolean) is
\[
  \sum_{v} w_p(v)\,\ell_{U,S,i}(v)
  \;=\;
  \begin{cases}
    p & \text{if } i \in S,\\
    0 & \text{if } i \in U \setminus S,\\
    1 & \text{otherwise.}
  \end{cases}
\]
\end{lemma}

\begin{lemma}[Averaging a product of two local factors]
\lean{RestrictionFourier.sum_varWeight_localFactor_mul}
\label{RestrictionFourier.sum_varWeight_localFactor_mul}
\leanok
\uses{BooleanAnalysis.boolToSign, LMN.varWeight, RestrictionFourier.localFactor}
\difficulty{2}
Fix a real parameter $p$, subsets $U, V, S \subseteq [n]$, and a coordinate $i \in [n]$.
Averaging over the three per-coordinate values $v$ — the free value and the two fixed
signs — with the per-variable weight $w_p(v)$, the product of the two local factors, one
for the pair $(U,S)$ and one for the pair $(V,S)$, satisfies
\[
  \sum_{v} w_p(v)\,\ell_{U,S}(i,v)\,\ell_{V,S}(i,v)
  = \begin{cases}
      p & i \in S,\\[2pt]
      1-p & i \notin S,\ i \in U,\ \text{and}\ i \in V,\\[2pt]
      0 & i \notin S\ \text{and exactly one of}\ i \in U,\ i \in V\ \text{holds},\\[2pt]
      1 & i \notin S,\ i \notin U,\ \text{and}\ i \notin V,
    \end{cases}
\]
where $w_p$ denotes the per-variable weight and $\ell_{W,S}$ the local factor determined
by the sets $W$ and $S$.
\end{lemma}

\begin{lemma}[Product of a three-valued weight over coordinates]
\lean{RestrictionFourier.prod_if_subset}
\label{RestrictionFourier.prod_if_subset}
\leanok
\difficulty{4}
Let $n$ be a natural number, let $p \in \bbr$, and let $S, U \subseteq \{1, \dots, n\}$.
For each coordinate $i$, assign the weight $p$ if $i \in S$, the weight $1 - p$ if $i
\in U \setminus S$, and the weight $1$ otherwise. Then the product of these weights over
all $n$ coordinates equals
\[
  p^{\abs{S}}\,(1-p)^{\abs{U \setminus S}}.
\]
\end{lemma}

\begin{theorem}[Expected Fourier coefficient under a random restriction]
\lean{RestrictionFourier.expectation_fourierCoeff_restrictBF}
\label{RestrictionFourier.expectation_fourierCoeff_restrictBF}
\leanok
\statementsource{boolean-ch04-dnf-switching-lmn}{
  pdf-fc447f326910-p006-b005,
  pdf-fc447f326910-p007-b001
}
\statementsource{boolean-ch03-spectral-learning-restrictions}{pdf-20e15a3d3d5e-p009-b001}
\uses{BooleanAnalysis.BooleanFunc, BooleanAnalysis.fourierCoeff, RestrictionFourier.fourierCoeff_restrictBF, RestrictionFourier.indicator_signProd_eq_prod, RestrictionFourier.sum_bernoulli_prod, RestrictionFourier.sum_varWeight_localFactor, SwitchingLemma2.Restriction, SwitchingLemma2.Restriction.freeVars, SwitchingLemma2.bernoulliRestrWeight}
\difficulty{5}
Let $f:\{0,1\}^n\to\bbr$ be a Boolean function, let $S\subseteq[n]$ be a frequency set,
and let $p\in\bbr$. For a restriction $\rho$ on $n$ variables, write $f_\rho$ for the
restricted function and $k(\rho)$ for the number of free variables of $\rho$. Summing
the Fourier coefficient $\widehat{f_\rho}(S)$ over all restrictions $\rho$, each
weighted by its Bernoulli weight $p^{k(\rho)}\bigl(\tfrac{1-p}{2}\bigr)^{\,n-k(\rho)}$,
gives
\[
  \sum_{\rho} p^{k(\rho)}\left(\tfrac{1-p}{2}\right)^{\,n-k(\rho)}\widehat{f_\rho}(S)
  \;=\; p^{\abs{S}}\,\hat f(S).
\]
(When $0\le p\le 1$ these weights form the distribution of a Bernoulli$(p)$-random
restriction, and the left-hand side is the expectation $\E_\rho[\widehat{f_\rho}(S)]$.)
\end{theorem}

\begin{theorem}[Expected squared Fourier coefficient under a random restriction]
\lean{RestrictionFourier.expectation_fourierCoeff_sq_restrictBF}
\label{RestrictionFourier.expectation_fourierCoeff_sq_restrictBF}
\leanok
\statementsource{boolean-ch04-dnf-switching-lmn}{
  pdf-fc447f326910-p006-b005,
  pdf-fc447f326910-p007-b001
}
\statementsource{boolean-ch03-spectral-learning-restrictions}{pdf-20e15a3d3d5e-p009-b001}
\uses{BooleanAnalysis.BooleanFunc, BooleanAnalysis.fourierCoeff, RestrictionFourier.fourierCoeff_restrictBF, RestrictionFourier.indicator_signProd_eq_prod, RestrictionFourier.localFactor, RestrictionFourier.prod_if_subset, RestrictionFourier.sum_bernoulli_prod, RestrictionFourier.sum_varWeight_localFactor_mul, SwitchingLemma2.Restriction, SwitchingLemma2.Restriction.freeVars, SwitchingLemma2.bernoulliRestrWeight}
\difficulty{6}
Fix a dimension $n$, a real number $p$, a Boolean function $f:\{0,1\}^n\to\bbr$, and a
set $S\subseteq[n]$. For a restriction $\rho$ on the $n$ variables, let $k(\rho)$ denote
its number of free variables and give it the Bernoulli weight
$w_p(\rho)=p^{\,k(\rho)}\bigl(\tfrac{1-p}{2}\bigr)^{\,n-k(\rho)}$; write $f_\rho$ for
the restriction of $f$ by $\rho$ and $\hat g(T)$ for the Fourier--Walsh coefficient of a
Boolean function $g$ at a set $T$. Then, with the left sum over all restrictions on the
$n$ variables and the right sum over all $U\subseteq[n]$,
\[
  \sum_{\rho} w_p(\rho)\,\widehat{f_\rho}(S)^2
  \;=\;
  \sum_{U\subseteq[n]}
  \begin{cases}
p^{\,\abs{S}}\,(1-p)^{\abs{U\setminus S}}\,\hat f(U)^2 & \text{if } S\subseteq U,\\[2pt]
    0 & \text{otherwise.}
  \end{cases}
\]
When $0\le p\le 1$ the weights $w_p(\rho)$ form a probability distribution, so the
left-hand side is the expectation of $\widehat{f_\rho}(S)^2$ over a Bernoulli$(p)$
random restriction.
\end{theorem}

\begin{lemma}[The sign product squares to one]
\lean{RestrictionFourier.signProd_sq}
\label{RestrictionFourier.signProd_sq}
\leanok
\uses{BooleanAnalysis.boolToSign_sq, SwitchingLemma2.Restriction}
\difficulty{2}
Let $\rho$ be a restriction on $n$ variables and let $T \subseteq [n]$, and write
$s(\rho, T) \in \bbr$ for the sign product of the bits that $\rho$ fixes on $T$, that
is, the product over $i \in T$ of the $\pm 1$-encoding of $\rho$'s value at $i$. Being a
product of values in $\{-1, 1\}$, it satisfies \[ s(\rho, T)^2 = 1. \]
\end{lemma}

\begin{theorem}[Distribution of a fixed coordinate set under a random restriction]
\lean{RestrictionFourier.bernoulliRestrProb_inter_freeVars}
\label{RestrictionFourier.bernoulliRestrProb_inter_freeVars}
\leanok
\uses{BooleanAnalysis.chiS, BooleanAnalysis.fourierCoeff, BooleanAnalysis.fourier_coeff_chi, RestrictionFourier.expectation_fourierCoeff_sq_restrictBF, RestrictionFourier.fourierCoeff_restrictBF, RestrictionFourier.signProd_sq, SwitchingLemma2.Restriction, SwitchingLemma2.Restriction.freeVars, SwitchingLemma2.bernoulliRestrProb}
\difficulty{5}
Fix $n$, a real parameter $p$, and two subsets $U, S \subseteq [n]$. Let $\rho$ be a
Bernoulli($p$) random restriction on $n$ variables, meaning each restriction is drawn
with weight $p^{\abs{J}}\bigl(\tfrac{1-p}{2}\bigr)^{\,n-\abs{J}}$, where $J$ denotes its
set of free variables. Then the probability that the free variables lying inside $U$ are
exactly those of $S$ is
\[
  \Pr_\rho\bigl[\,U \cap J = S\,\bigr] \;=\;
  \begin{cases}
    p^{\abs{S}}\,(1-p)^{\abs{U \setminus S}}, & S \subseteq U,\\[2pt]
    0, & \text{otherwise.}
  \end{cases}
\]
\end{theorem}

%%%%%%%%%%%%%%%%%%%%%%%%%%%%%%%%%%%%%%%%%%%%%%%%%%%%%%%%%%%%%%%%%%%%%%%%%%%%%%%
%%% added by the blueprint pipeline
\section{Additional declarations}

\begin{definition}[Restriction of a real-valued Boolean function]
\lean{RestrictionFourier.restrictBF}
\statementsource{boolean-ch03-spectral-learning-restrictions}{pdf-20e15a3d3d5e-p007-b001}
\label{RestrictionFourier.restrictBF}
\leanok
\statementsource{boolean-ch04-dnf-switching-lmn}{pdf-fc447f326910-p006-b004}
\uses{BooleanAnalysis.BooleanFunc, SwitchingLemma2.Restriction, SwitchingLemma2.Restriction.extend}
Given $f : \{0,1\}^n \to \bbr$ and a restriction $\rho$, the restricted function
$f_\rho = \mathrm{restrictBF}\,f\,\rho$ is defined by
\[
  (\mathrm{restrictBF}\,f\,\rho)(x) \;=\; f(\rho.\mathrm{extend}\,x),
\]
so free coordinates are read from the input $x$ and fixed coordinates from $\rho$. It is
the real-valued analogue of $\mathrm{restrictFn}$.
\end{definition}

\begin{definition}[Sign product of the fixed bits]
\lean{RestrictionFourier.signProd}
\label{RestrictionFourier.signProd}
\leanok
\uses{BooleanAnalysis.boolToSign, SwitchingLemma2.Restriction}
For a restriction $\rho$ and a set $T \subseteq [n]$,
\[
  \mathrm{signProd}\,\rho\,T \;=\; \prod_{i \in T} \mathrm{boolToSign}\bigl((\rho\,i).\mathrm{getD}\ \mathtt{false}\bigr)
  \;\in\; \bbr,
\]
the product of the $\pm1$-encodings of the bits that $\rho$ fixes on $T$. On coordinates
where $\rho$ is free the default value $\mathtt{false}$ is used; this junk value is
harmless because the definition is only ever paired with indicators forcing $T$ to consist
of fixed coordinates.
\end{definition}
